Artificial autonomy requires rethinking the role of reward in Reinforcement Learning. In traditional approaches, the agent learns from external rewards defined by the designer, which may produce efficient behaviors but limits the emergence of more autonomous forms of action. Inspired by Enactive Artificial Intelligence, this work investigates agents whose actions are guided by internal viability conditions. To this end, the \textit{Health-Gathering} experimental task from the VizDoom platform is used, in which an agent continuously loses health in a room with an acidic floor and perceives the environment through vision. Instead of explicitly rewarding the collection of \textit{medikits}, items that restore the agent's health, an artificial form of bodily precariousness is introduced. As the agent remains without collecting \textit{medikits}, its vision progressively deteriorates and is restored when it collects them. The agent's motivation is produced by an intrinsic curiosity-based reward calculated using the \textit{Random Network Distillation} (RND) method. Since curiosity depends on visual observations, perceptual degradation reduces the agent's ability to obtain intrinsic reward. Thus, remaining viable becomes a condition for continuing to perceive, explore, and act. The learned policy will be analyzed using visual interpretability techniques to understand how precariousness, perception, and curiosity may promote adaptive behaviors without directly imposing a solution through extrinsic rewards. Therefore, if autonomy arises from precariousness, then being intelligent fundamentally means being engaged in a continuous process of negotiating with the possibility of ceasing to exist.

% Separate keywords with periods
\keywords{Autonomous agents. Enactive Artificial Intelligence. Curiosity. Deep Reinforcement Learning. Precariousness. Intrinsic motivation.}